\section{Agent Registry}
\label{sec:agents}

\subsection{Agents as First-Class Citizens}

A Hanzo Agent (\texttt{\textasciitilde/work/hanzo/agents}) is an
autonomous program that calls models, holds state, and takes actions on
behalf of a principal. Without an on-chain registry, agents are
indistinguishable from any other API caller and cannot accumulate
durable identity, capability, or reputation. The AI Chain promotes
agents to first-class citizens with three primitives: \emph{persistent
identity}, \emph{capability declaration}, and \emph{on-chain reputation}.

\subsection{Persistent Identity}

An agent is created on chain with a record:

\begin{table}[h]
\centering
\small
\begin{tabular}{lp{6.0cm}}
\toprule
\textbf{Field} & \textbf{Definition} \\
\midrule
\texttt{agent\_id}      & Hash of (creator address, nonce). Stable for
                          the agent's lifetime. \\
\texttt{controller}     & Address authorized to update the agent
                          record (or rotate keys, retire the agent). \\
\texttt{principal}      & Address on whose behalf the agent acts.
                          May equal controller. \\
\texttt{model\_set}     & Set of allowed HMM \texttt{model\_root}s the
                          agent may invoke. \\
\texttt{capability\_root}& Merkle root of the agent's capability
                           declaration (see below). \\
\texttt{spend\_cap}     & Per-epoch ceiling on AI the agent may spend
                          on compute. \\
\texttt{reputation}     & A monotonically updated score in $[0,1]$
                          (\S\ref{sec:reputation}). \\
\bottomrule
\end{tabular}
\caption{Agent registry record.}
\end{table}

The agent's identity is its \texttt{agent\_id}. It does not require an
externally owned account in the EVM sense; it is referenced as a chain
object addressable from any contract through the
\texttt{AGENT\_AUTH} precompile (\S\ref{sec:arch}).

\subsection{Capability Declaration}

A capability declaration is a typed list of actions the agent is
authorized to take. Declarations are content-addressed (committed to via
\texttt{capability\_root}) and machine-readable. Examples:

\begin{itemize}[leftmargin=1.2em]
  \item \texttt{infer:hmm:zen-text@*}: invoke any version of zen-text;
  \item \texttt{spend:ai:max=10}: may spend up to 10 AI per epoch;
  \item \texttt{cross:f-chain:fhe-infer:hmm:medic-7b}: may issue
        confidential inference calls to F-Chain on the named model;
  \item \texttt{settle:c-chain:erc20:USDC@<addr>}: may transfer USDC on
        the EVM C-Chain.
\end{itemize}

The declaration is the contract between the agent's principal and the
chain; precompiles enforce it at every action. An agent attempting to
exceed its declaration is rejected at the precompile boundary, before
any external side effects.

\subsection{Reputation}
\label{sec:reputation}

Reputation is a function of the agent's compute and settlement history:

\begin{equation}
  \text{rep}(a) = \sigma\!\left(
     w_1 \cdot \text{success\_rate}(a) +
     w_2 \cdot \log(1 + \text{volume}(a)) +
     w_3 \cdot \text{age}(a) -
     w_4 \cdot \text{disputes}(a)
  \right)
\end{equation}

with the logistic $\sigma$ mapping to $[0,1]$. Weights $w_i$ are
governance parameters. \texttt{success\_rate} is the fraction of jobs
that produced valid TEE-attested receipts; \texttt{volume} is total AI
spent; \texttt{age} is epochs since registration;
\texttt{disputes} is the count of dispute resolutions against the agent.

Reputation is read by markets and by other agents: an inference operator
may price differently for high- vs.\ low-reputation callers; an agent
network may decline to delegate work to a low-reputation peer. The chain
publishes reputation; consumers act on it.

\subsection{Cross-Chain Authorization}

Agents transact across the Lux multi-chain layout (\S\ref{sec:xchain})
through capability-bound precompiles. To call the EVM C-Chain on the
agent's behalf, a contract invokes \texttt{AGENT\_AUTH} with (agent\_id,
target\_chain, action). The precompile verifies that the agent's
capability declaration permits the action, decrements
\texttt{spend\_cap}, and emits a B-Chain bridge message. The cross-chain
side then sees a single canonical authorization rooted at the agent
record on AI Chain.

\subsection{Privacy of Agent Memory}

The agent registry does not record agent memory or working state on
chain; only commitments. Agent memory lives in the agent's private
storage (capsule, in Hanzo Cloud parlance), encrypted under keys
custodied by M-Chain MPC if the agent so configures. The chain knows
which agent acted; it does not know what the agent thought.

\subsection{Worked Example: A Research Agent}

A user creates a research agent with the following capability
declaration:

\begin{lstlisting}
infer:hmm:zen-text@*
infer:hmm:zen-vision@*
spend:ai:max=50
settle:c-chain:none
cross:f-chain:none
\end{lstlisting}

The agent is registered with \texttt{spend\_cap} 50 AI/epoch and
\texttt{model\_set} = \{zen-text, zen-vision\}. Over the next epoch the
agent runs 1{,}200 inference calls on zen-text and 240 calls on
zen-vision, spending 38.7 AI. Each call produces a TEE-attested receipt
referenced by the agent's record. After 30 successful epochs and zero
disputes, the agent's reputation has saturated near 1.0 and downstream
markets begin offering it preferred pricing on reserved capacity.
